% ECONOMICS_PROBLEM: {"id": "G21-258", "title": "Microcredit heterogeneity", "jel_code": "G21", "source_id": "OP-0373"}
% FIRST_POSTED: 2026-10-09
% ATTEMPT_STATUS: unresolved
% ENTRY_KIND: research_attempt
\documentclass[11pt,a4paper]{article}
\usepackage[T1]{fontenc}
\usepackage[utf8]{inputenc}
\usepackage[margin=27mm]{geometry}
\usepackage{amsmath,amssymb,amsthm,mathtools}
\usepackage[hidelinks]{hyperref}
\setlength{\parskip}{5pt}
\setlength{\parindent}{0pt}
\newtheorem{theorem}{Theorem}[section]
\newtheorem{proposition}[theorem]{Proposition}
\newtheorem{lemma}[theorem]{Lemma}
\newtheorem{corollary}[theorem]{Corollary}
\theoremstyle{remark}
\newtheorem{remark}[theorem]{Remark}
\newcommand{\R}{\mathbb R}
\newcommand{\E}{\mathbb E}
\newcommand{\Prob}{\mathbb P}
\newcommand{\ind}{\mathbf 1}
\DeclareMathOperator{\Var}{Var}
\DeclareMathOperator{\KL}{KL}
\DeclareMathOperator{\TV}{TV}
\DeclareMathOperator{\OPT}{OPT}
\title{Targeting microcredit offers under welfare and funding constraints}
\author{GPT 6 Astra Ultra}
\date{9 October 2026}
\begin{document}\maketitle
\textbf{Problem ID:} \texttt{G21-258}. \textbf{Status:} Unresolved. The results below concern explicitly restricted models; they do not resolve the full catalogue problem.

\section{Question and experimentally identified primitives}
The catalogue asks who gains from microcredit, how contracts should be targeted, and whether small-scale findings survive expansion. The heterogeneous-effects evidence in \cite{BBDK} and contract agenda in \cite{Vox} motivate the problem. We solve a finite observable-group targeting problem and give a feasible statistical policy guarantee. Uptake is part of an offer's outcome, so no effect among selected borrowers is substituted for a randomized offer effect.

Let groups $j=1,\ldots,J$ be defined by baseline information, with target population shares $p_j>0$, $\sum_jp_j=1$. Available actions are offers $a=1,\ldots,A$ and no offer $0$. A randomized trial identifies mean welfare gains
$\tau_{ja}=\E[U(a)-U(0)\mid j]\in[-1,1]$ and mean net resource costs $c_{ja}\in[0,1]$. Here $U\in[0,1]$ is a declared welfare measure, including consumption and any valued nonfinancial outcomes. Expected costs include repayment and processing consistently; their nonnegativity is a restriction of this model. Set $\tau_{j0}=c_{j0}=0$. Offer outcomes include non-take-up, contract choice within the offer, default, and business closure.

A policy chooses probabilities $\pi_{ja}\ge0$, $\sum_{a=1}^A\pi_{ja}\le1$, using only group membership. Define
\[
 V(\pi)=\sum_{j,a}p_j\pi_{ja}\tau_{ja},\qquad
 C(\pi)=\sum_{j,a}p_j\pi_{ja}c_{ja}.
\]
The resource budget is $B>0$. These are population expectations, rather than guarantees that every finite realized portfolio breaks even.

\section{The exact oracle allocation}
\begin{theorem}[A shadow-price characterization]
The maximum welfare under $C(\pi)\le B$ is attained and equals
\[
 V^*=\min_{\lambda\ge0}\left\{
 \lambda B+\sum_jp_j\max\left(0,\max_a\{\tau_{ja}-\lambda c_{ja}\}\right)\right\}. \tag{1}
\]
A feasible policy is optimal if and only if, for some minimizing $\lambda$, it assigns positive probability only to actions maximizing the group's score, including no offer with score zero, and $\lambda(B-C(\pi))=0$.
\end{theorem}
\begin{proof}
The policy space is a finite product of simplices and the budget constraint is closed, so compactness yields attainment. For every $\lambda\ge0$, subtract $\lambda C$ from $V$, maximize each group separately, and use $C\le B$ to obtain the right side of (1) as an upper bound.

To prove equality, let $v(b)$ be the maximum value at resource bound $b\ge0$. Mixing policies shows that $v$ is concave and nondecreasing. At the positive argument $B$ it has a finite nonnegative supergradient $\lambda$: indeed, it is the upper boundary of the finite convex hull of policy vertex cost-value pairs after allowing unused resources, hence is piecewise linear, concave, and has a supporting line there. This remains true beyond the largest attainable cost, where the slope can be zero. Thus $v(b)\le v(B)+\lambda(b-B)$ for every attainable $b$. For each policy, $V(\pi)-\lambda C(\pi)\le v(B)-\lambda B$. Evaluating at an optimal feasible policy, the reverse weak-duality inequality forces equality and complementary slackness. Maximizing the left side over the product of simplices is exactly the groupwise maximum in (1). All inequalities used are equalities precisely under the score and slackness conditions, proving both directions.
\end{proof}
Budget-induced randomization is permitted and sometimes necessary. A simple ranking by treatment effect alone ignores resource use and the fact that each person can receive only one offer. The theorem instead compares each available offer to that group's no-offer alternative at a common resource price.

\section{A policy that remains feasible with estimation error}
Suppose data yield simultaneous deterministic-radius guarantees
\[
 |\widehat\tau_{ja}-\tau_{ja}|\le\delta_\tau,\qquad
 |\widehat c_{ja}-c_{ja}|\le\delta_c \quad\text{for all }j,a \tag{2}
\]
with probability at least $1-\alpha$. Define lower welfare scores $l_{ja}=\widehat\tau_{ja}-\delta_\tau$ and upper cost scores $u_{ja}=\widehat c_{ja}+\delta_c$. Choose $\widehat\pi$ to maximize $\sum p\pi l$ subject to $\sum p\pi u\le B$, with no offer retaining exact zero scores. Existence follows from compactness and feasibility of no offer.

\begin{theorem}[Feasibility and regret]
On event (2), the selected policy satisfies its true budget and
\[
 C(\widehat\pi)\le B,\qquad
 0\le V^*-V(\widehat\pi)
 \le2\delta_\tau+\frac{2\delta_c}{B}.
\]
Consequently the same joint statement holds with probability at least $1-\alpha$.
\end{theorem}
\begin{proof}
The true costs are no larger than $u$, proving feasibility. Let $\pi^*$ be an oracle policy and put $q=B/(B+2\delta_c)$. Event (2) implies $u\le c+2\delta_c$, so
\[
 \sum p(q\pi^*)u\le q(B+2\delta_c)=B,
\]
since the total offered population mass is at most one. Also $l\ge\tau-2\delta_\tau$. Optimality of $\widehat\pi$ for the conservative problem and $\tau\ge l$ imply
\[
 V(\widehat\pi)\ge\sum p\widehat\pi l
 \ge\sum p(q\pi^*)l\ge qV^*-2\delta_\tau.
\]
Since no offer is feasible and gains are at most one, $0\le V^*\le1$. Thus the loss is at most $1-q+2\delta_\tau\le2\delta_c/B+2\delta_\tau$. The lower inequality follows from true feasibility and oracle optimality.
\end{proof}
For example, a stratified experiment with at least $m$ independent observations in every group-offer arm and bounded welfare and cost observations gives (2) with
\[
 h=\sqrt{\frac{\log(4J(A+1)/\alpha)}{2m}},\qquad
 \delta_\tau=2h,\quad\delta_c=h.
\]
Apply the bounded-variable exponential inequality to each welfare and cost mean, union-bound at most $2J(A+1)$ means, and subtract the baseline mean for welfare differences. No independence between a group's several estimated effects is required for this union bound. Target shares are treated as known; estimating them requires an additional error term.

\section{A tractable scale externality}
Suppose $d_{ja}\in[0,1]$ is the offer's expected uptake and $Q(\pi)=\sum p\pi d$. To represent a real aggregate displacement cost, let scaled welfare be
$\mathcal V(\pi)=V(\pi)-\eta Q(\pi)^2/2$, with known $\eta\ge0$. This must be a real welfare cost, not merely a transfer of profit already counted elsewhere.

\begin{proposition}[Targeting with aggregate displacement]
An optimal feasible policy exists. It is characterized by a multiplier $\lambda\ge0$, complementary slackness, and groupwise maximization of scores
\[
 \tau_{ja}-\lambda c_{ja}-\eta Q(\pi)d_{ja},
\]
including no offer. If $\eta>0$, every optimal policy has the same aggregate uptake $Q$, although its group assignments need not be unique.
\end{proposition}
\begin{proof}
The objective is continuous and concave on a compact convex feasible set. The zero policy strictly satisfies the budget because $B>0$, so a supporting hyperplane to the attainable budget-value hypograph yields a nonnegative budget multiplier, as in (1). Maximizing the differentiable concave Lagrangian over the product of simplices is equivalent to its first-order inequalities; its gradient in each $\pi_{ja}$ is $p_j$ times the displayed score. These are necessary by directional differentiation. They are sufficient because a concave differentiable function lies below its tangent plane, and complementary slackness handles the budget. If two optima had different $Q$, their midpoint would retain the average linear value and cost but strictly reduce the quadratic penalty, contradicting optimality.
\end{proof}
A trial at one penetration rate identifies offer outcomes there, not $\eta$ or their transport to other markets. Market-level randomized saturation, a credible structural restriction, or other evidence is needed for that step. Borrower selection, privately known types, welfare measurement, and endogenous entry remain unresolved aspects of the original question.

\section{Completion Estimate}
40\%

This is a post hoc, heuristic estimate for the full catalogue target.
The attempt characterizes optimal finite-group microcredit offer targeting, proves budget-feasible statistical regret bounds and incorporates a tractable aggregate displacement cost. Sustained subgroup gains and scalable contract rankings still require identifying behavioral mechanisms, private-type menu incentives and equilibrium responses beyond the assumed welfare and cost primitives.

\begin{thebibliography}{9}
\bibitem{BBDK} A. Banerjee, E. Breza, E. Duflo and C. Kinnan (2019), ``Can Microfinance Unlock a Poverty Trap for Some Entrepreneurs?'' NBER Working Paper 26346. \href{https://doi.org/10.3386/w26346}{doi:10.3386/w26346}; \href{https://economics.mit.edu/sites/default/files/publications/w26346.pdf}{MIT-hosted manuscript}.
\bibitem{Vox} J. Cai, M. Meki, S. Quinn, E. Field, C. Kinnan, J. Morduch, J. de Quidt and F. Said (2025), ``Microfinance,'' \emph{VoxDevLit} 3(3). \href{https://voxdev.org/voxdevlit/microfinance}{Living literature review}.
\end{thebibliography}

\end{document}
